\bookchapter{Κεφάλαιο 67: Γυάλισ$\mu$α και κυκλοφορία — δοκι$\mu$ές, debugging και ανέβασ$\mu$α στο App Store}

\chapterintro{Έφτιαξες την εφαρ$\mu$ογή — και $\mu$ετά; Σε αυτό το κεφάλαιο απαντιούνται τα τρία ζητή$\mu$ατα της πορείας «από το να τρέχει» στο «να το παραδίδεις»: πώς κάνεις αποδοτικό debugging για να βρίσκεις bugs, πώς γράφεις δοκι$\mu$ές για να $\mu$ην επιστρέφουν παλιές βλάβες, και πώς ανεβάζεις την εφαρ$\mu$ογή στο App Store για να τη κατεβάζει όλος ο κόσ$\mu$ος. Το ανέβασ$\mu$α δεν είναι απλώς τελετουργικό — είναι το ορόση$\mu$ο που $\mu$ετατρέπει ένα έργο σε προϊόν.}

\booksection{Τα τρία βασικά εργαλεία debugging: print, breakpoints, preview}

\begin{center}
\begin{tabularx}{\linewidth}{|X|X|X|}
\hline
\textbf{Εργαλείο} & \textbf{Τι κάνει} & \textbf{Πώς το χρησι$\mu$οποιείς} \tabularnewline
\hline
print + κονσόλα & Ε$\mu$φάνιση logs εκτέλεσης & Στον κώδικα γράφεις print("τι$\mu$ή=\textbackslash{}(x)") και βλέπεις το αποτέλεσ$\mu$α στην κονσόλα κάτω στο Xcode \tabularnewline
Debugging $\mu$ε breakpoints & Στα$\mu$ατάς το πρόγρα$\mu$$\mu$α σε $\mu$ια γρα$\mu$$\mu$ή και βλέπεις τις τι$\mu$ές των $\mu$εταβλητών & Ένα κλικ δίπλα στον αριθ$\mu$ό γρα$\mu$$\mu$ής βάζει $\mu$πλε ση$\mu$άδι· η εκτέλεση στα$\mu$ατά εκεί όταν φτάσει, και προχωράς βή$\mu$α-βή$\mu$α \tabularnewline
Προεπισκόπηση SwiftUI + debugger & Βλέπεις τη διεπαφή και την ιεραρχία των views σε πραγ$\mu$ατικό χρόνο & Κα$\mu$βάς \#Preview + επιθεωρητής View Hierarchy κατά την εκτέλεση \tabularnewline
\hline
\end{tabularx}
\end{center}

\begin{center}\small\color{gray}Τα τρία βασικά εργαλεία debugging\end{center}

\begin{tcbinfo}{Μικρή συ$\mu$βουλή: Πώς διαβάζεις ένα crash;}
Όταν ένα πρόγρα$\mu$$\mu$α Swift κρασάρει, η κονσόλα τυπώνει «Fatal error» $\mu$αζί $\mu$ε την αιτία, και το Xcode στα$\mu$ατά στη γρα$\mu$$\mu$ή του σφάλ$\mu$ατος και τη τονίζει. Τα ιστορικά crashes σε πραγ$\mu$ατική συσκευή $\mu$πορείς να τα δεις στο Xcode Window Organizer Crashes (θέλει πληρω$\mu$ένο λογαριασ$\mu$ό). Στην αναζήτηση βάλε προτεραιότητα στο «πρώτη φράση της αιτίας του σφάλ$\mu$ατος + swiftui» — $\mu$όνο ακριβής ερώτηση δίνει ακριβή απάντηση.

\end{tcbinfo}

\booksection{Δοκι$\mu$ές: ασφάλεια για τον κώδικά σου}

Οι χειροκίνητες δοκι$\mu$ές δεν αξίζουν πολλή ε$\mu$πιστοσύνη — αλλάζεις αυτή τη λειτουργία και $\mu$ια άλλη χαλάει διακριτικά. Οι δοκι$\mu$ές είναι «φάρ$\mu$ακο της $\mu$ετα$\mu$έλειας» και ταυτόχρονα «αλεξίσφαιρο γιλέκο». Το iOS χρησι$\mu$οποιεί το framework XCTest, $\mu$ε τα εξής είδη:

\begin{center}
\begin{tabularx}{\linewidth}{|X|X|X|X|}
\hline
\textbf{Επίπεδο} & \textbf{Τι δοκι$\mu$άζει} & \textbf{Target} & \textbf{Ταχύτητα} \tabularnewline
\hline
Δοκι$\mu$ή $\mu$ονάδας (unit test) & Καθαρή λογική (υπολογισ$\mu$οί, έλεγχοι) & MyFirstAppTests & Δευτερόλεπτα \tabularnewline
Δοκι$\mu$ή UI (UI test) & Πραγ$\mu$ατική συ$\mu$περιφορά της διεπαφής & MyFirstAppUITests & Λεπτά \tabularnewline
Χειροκίνητη δοκι$\mu$ή & Αίσθηση και ε$\mu$πειρία από πραγ$\mu$ατικό άνθρωπο & —— & Χρειάζεται πάντα \tabularnewline
\hline
\end{tabularx}
\end{center}

\begin{center}\small\color{gray}Επίπεδα δοκι$\mu$ών στο iOS (το template του έργου περιλα$\mu$βάνει ήδη δύο Test Target)\end{center}

\begin{lstlisting}[language=Swift]
import XCTest

// Διαχωρίζου$\mu$ε τη λογική "$\mu$ετατροπή κει$\mu$ένου ποσού σε αριθ$\mu$ό" σε καθαρή συνάρτηση
func parseAmount(_ input: String) -> Double? {
    guard let v = Double(input.trimmingCharacters(in: .whitespaces)),
          v > 0 else { return nil }
    return v
}

final class AmountTests: XCTestCase {
    func testEmptyInputReturnsNil() {
        XCTAssertNil(parseAmount(""))
    }
    func testNegativeReturnsNil() {
        XCTAssertNil(parseAmount("-5"))
    }
    func testNormalValueParses() {
        XCTAssertEqual(parseAmount("12.5"), 12.5)
    }
}
\end{lstlisting}

Τρέξι$\mu$ο δοκι$\mu$ών: CmdU. Όταν όλα είναι πράσινα, αλλάζεις κώδικα $\mu$ε ησυχία. Οι έλεγχοι ποσών και οι υπολογισ$\mu$οί συνόλων στο βιβλίο εξόδων πρέπει όλοι να καλύπτονται από unit tests.

\booksection{Υπογραφή και πακετάρισ$\mu$α: βάλε τη «σφραγίδα» στην εφαρ$\mu$ογή}

Κάθε εφαρ$\mu$ογή iOS πρέπει να είναι υπογεγρα$\mu$$\mu$ένη για να εγκαθίσταται — αποδεικνύει ότι «το πακέτο προέρχεται από σένα και δεν έχει παραποιηθεί». Η ιεραρχική σχέση:

\begin{itemize}
  \item Στάδιο $\mu$άθησης: δωρεάν Apple ID + αυτό$\mu$ατη υπογραφή στο Xcode (Personal Team), τρέχει σε πραγ$\mu$ατική συσκευή, $\mu$ε ανανέωση κάθε 7 η$\mu$έρες·
  \item Στάδιο διανο$\mu$ής: πληρω$\mu$ένος λογαριασ$\mu$ός προγρα$\mu$$\mu$ατιστή (99 δολάρια/έτος) που εκδίδει πιστοποιητικά και provisioning profiles, για να $\mu$πορείς να ανεβάσεις στο κατάστη$\mu$α και να εγκαθιστάς $\mu$ακροπρόθεσ$\mu$α.
\end{itemize}

Πακετάρισ$\mu$α: στο $\mu$ενού του Xcode, Product Archive (αρχειοθέτηση)· $\mu$ετά την ολοκλήρωση του build, στο Organizer διάλεξε Distribute App App Store Connect ανέβασ$\mu$α. Η αρχειοθέτηση πρέπει να γίνει $\mu$ε στόχο το Generic iOS Device (όχι τον εξο$\mu$οιωτή) — βή$\mu$α που οι αρχάριοι ξεχνάνε συχνά.

\textbf{Προσοχή στην παγίδα: Χάθηκε το πιστοποιητικό = στα$\mu$ατάνε οι ενη$\mu$ερώσεις}\par

Αν χάσεις το πιστοποιητικό υπογραφής και το ιδιωτικό κλειδί, δεν $\mu$πορείς πια να στείλεις ενη$\mu$ερώσεις στην παλιά εφαρ$\mu$ογή — οι χρήστες $\mu$πορούν να εγκαταστήσουν $\mu$όνο «ολοκαίνουργια» εφαρ$\mu$ογή. Κράτα προσεκτικά backup του πιστοποιητικού (εξαγωγή .p12 από το Keychain) και ση$\mu$είωσε καλά τον κωδικό. Τυχερά, η Apple προσφέρει «αυτό$\mu$ατη υπογραφή» και διαχείριση πιστοποιητικών στο cloud (τσέκαρε στο Xcode το Automatically manage signing), που σου γλιτώνει τις περισσότερες χειροκίνητες παγίδες.

\booksection{Η πλήρης διαδικασία ανεβάσ$\mu$ατος στο App Store}

Εγγραφή λογαριασ$\mu$ού developer

Προετοι$\mu$ασία υλικών για το κατάστη$\mu$α

Ανέβασ$\mu$α πακέτου build

Συ$\mu$πλήρωσε στοιχεία ιδιωτικότητας και ελέγχου

Έλεγχος Apple (1-7 $\mu$έρες)

Κυκλοφορία

\begin{center}\small\color{gray}Εικόνα 5: Έξι βή$\mu$ατα για το App Store\end{center}

\begin{itemize}
  \item Εγγραφή: developer.apple.com, προσωπικός λογαριασ$\mu$ός 99 δολάρια/έτος·
  \item Υλικό: εικονίδιο (1024×1024, χωρίς στρογγυλε$\mu$ένες γωνίες — η Apple κόβει $\mu$όνη της), στιγ$\mu$ιότυπα (για όλα τα $\mu$εγέθη iPhone), περιγραφή, σύνδεσ$\mu$ος πολιτικής απορρήτου·
  \item Ερωτη$\mu$ατολόγιο απορρήτου: δήλωσε ειλικρινά ποια δεδο$\mu$ένα συλλέγεις — η εφαρ$\mu$ογή εξόδων χωρίς σύνδεση internet $\mu$πορεί να δηλώσει «δεν συλλέγω κανένα δεδο$\mu$ένο»·
  \item Έλεγχος (review): η πρώτη υποβολή συχνά περι$\mu$ένει 1-7 η$\mu$έρες· αν απορριφθεί, διόρθωσε ση$\mu$είο-ση$\mu$είο σύ$\mu$φωνα $\mu$ε τις «ση$\mu$ειώσεις του ελέγχου» — οι λόγοι απόρριψης της Apple συνήθως είναι γρα$\mu$$\mu$ένοι πολύ συγκεκρι$\mu$ένα, απλώς ακολούθησέ τους·
  \item TestFlight: πριν την κυκλοφορία, στείλε τη δοκι$\mu$αστική έκδοση σε φίλους για εσωτερική δοκι$\mu$ή (έως 10.000 άτο$\mu$α), $\mu$άζεψε έναν γύρο σχόλια και $\mu$ετά κάνε την επίση$\mu$η υποβολή — απόλυτα επαγγελ$\mu$ατική εικόνα.
\end{itemize}

\booksection{Πώς φτιάχνεις το υλικό του καταστή$\mu$ατος: η πρώτη εντύπωση είναι οι λήψεις}

Η σελίδα του καταστή$\mu$ατος είναι η βιτρίνα που βλέπει ο χρήστης «όταν κάνει βόλτα στο κατάστη$\mu$α». Τρία στοιχεία καθορίζουν την πρώτη εντύπωση:

\begin{itemize}
  \item Εικονίδιο: 1024×1024, $\mu$ην «συ$\mu$βιβαστείς» $\mu$ε ένα στιγ$\mu$ιότυπο. Αν δεν έχεις σχεδιαστικές δεξιότητες, φτιάξε $\mu$ε εργαλείο AI ένα απλό, καθαρό σχή$\mu$α — είναι αποδεκτό $\mu$όνο αν το αναγνωρίζεις $\mu$ε $\mu$ια $\mu$ατιά ακό$\mu$α και συρρικνω$\mu$ένο στην αρχική οθόνη·
  \item Στιγ$\mu$ιότυπα: τουλάχιστον 2 σετ (στις βασικές αναλύσεις iPhone), που δείχνουν τις πιο βασικές λειτουργίες· σε κάθε ένα πρόσθεσε $\mu$ια φράση απλή, ανθρώπινη («Καταγραφή εξόδου σε ένα δευτερόλεπτο» είναι δέκα φορές καλύτερη από «Αποδοτική λύση λογιστικής»)·
  \item Περιγραφή: η πρώτη φράση να λέει καθαρά «τι είναι + σε ποιον απευθύνεται»· $\mu$ην στοιβάζεις «το πιο επαγγελ$\mu$ατικό, το πιο δυνατό» — η ειλικρίνεια αποδίδει περισσότερο από το $\mu$εγαλόπνοο.
\end{itemize}

\begin{tcbpractice}{Φτιάξε το $\mu$όνο σου: «παράδωσε» την εφαρ$\mu$ογή εξόδων $\mu$ια φορά}
\booksection{Αναφορές crash και $\mu$ικρή λίστα βελτιστοποίησης απόδοσης}

Πριν ανέβει η εφαρ$\mu$ογή στο κατάστη$\mu$α, αξίζει $\mu$ισή ώρα για καθένα από τα δύο αυτά:

\booksubsection{Παρακολούθηση crashes}

Το ότι δεν κρασάρει στις δικές σου δοκι$\mu$ές δεν ση$\mu$αίνει ότι δεν θα κρασάρει στο κινητό του χρήστη. Η Apple παρέχει δωρεάν αναφορές crash στο Xcode Organizer ($\mu$ε πληρω$\mu$ένο λογαριασ$\mu$ό βλέπεις και τα crash logs από πραγ$\mu$ατικές συσκευές)· πιο ανώδυνη λύση είναι να ενσω$\mu$ατώσεις SDK τρίτου (π.χ. Firebase Crashlytics). Κυκλοφορία χωρίς παρακολούθηση crashes ισούται $\mu$ε οδήγηση $\mu$ε κλειστά $\mu$άτια.

\booksubsection{Τέσσερα ση$\mu$εία αυτοελέγχου απόδοσης}

\begin{enumerate}
  \item Λίστες $\mu$ε LazyVStack / List: $\mu$ακριές λίστες $\mu$ην τις στήνεις $\mu$ε απλό VStack — το lazy loading είναι το πρώτο $\mu$άθη$\mu$α απόδοσης·
  \item Μην πετάς $\mu$έσα πρωτότυπες εικόνες: χρησι$\mu$οποίησε AsyncImage (για εικόνες από το internet) και την cache του συστή$\mu$ατος — $\mu$ια πρωτότυπη εικόνα 5MB αρκεί για να «πνίξει» παλιές συσκευές·
  \item Μην κάνεις βαριές δουλειές στο κύριο νή$\mu$α (main thread): δίκτυο και $\mu$εγάλα αρχεία πάνε όλα $\mu$ε async/await·
  \item Μέτρησε $\mu$ε το Instruments: ο αναλυτής απόδοσης που έρχεται $\mu$ε το Xcode (Time Profiler, έλεγχος διαρροών $\mu$νή$\mu$ης) — πρώτα $\mu$έτρησε, $\mu$ετά βελτιστοποίησε, $\mu$ην δουλεύεις $\mu$ε το ένστικτο.
\end{enumerate}

(1) Πρόσθεσε 3 unit tests σε λογική όπως η parseAmount (να τρέχουν πράσινα $\mu$ε CmdU)· (2) $\mu$ε breakpoints πιάσε $\mu$όνος σου ένα crash που έφτιαξες εσύ και διόρθωσέ το· (3) $\mu$ε Product Archive φτιάξε ένα πακέτο και $\mu$έσω TestFlight κάλεσε φίλους για εσωτερική δοκι$\mu$ή· (4) αν έχεις περιθώριο, εγγράψου στο πρόγρα$\mu$$\mu$α προγρα$\mu$$\mu$ατιστών και ανέβασε πραγ$\mu$ατικά το βιβλίο εξόδων στο κατάστη$\mu$α. Τα τρία πρώτα βή$\mu$ατα $\mu$πορούν να τελειώσουν απόψε.

\end{tcbpractice}

\begin{tcbsummary}{Σύνοψη κεφαλαίου}
\begin{itemize}
  \item Τρία εργαλεία debugging: logs $\mu$ε print, breakpoints, προεπισκόπηση SwiftUI και επιθεωρητής views
  \item XCTest: unit tests (CmdU) + UI tests + χειροκίνητες δοκι$\mu$ές — το template τα έχει ήδη έτοι$\mu$α
  \item Υπογραφή: δωρεάν λογαριασ$\mu$ός $\mu$ε αυτό$\mu$ατη υπογραφή (ανανέωση κάθε 7 η$\mu$έρες)· $\mu$όνο πληρω$\mu$ένος λογαριασ$\mu$ός $\mu$πορεί να διανεί$\mu$ει
  \item Πακετάρισ$\mu$α: Generic iOS Device + Archive + Distribute· οπωσδήποτε κράτα backup των πιστοποιητικών
  \item Έξι βή$\mu$ατα κυκλοφορίας: εγγραφή υλικό ανέβασ$\mu$α ερωτη$\mu$ατολόγιο απορρήτου έλεγχος στο ράφι· πρώτα εσωτερική δοκι$\mu$ή $\mu$ε TestFlight
\end{itemize}

\end{tcbsummary}
